%%%%%%%%%%%%%%%%%%%%%%%%%%%%%%%%%%%%%%%%%%%%%%%%%%%%%%%%%%%%%%%%%%%%%%%%%%%%%%%
\documentclass[hyperref={pdfpagelabels=false},compress,table]{beamer} % 在Mac下无法编译
% \documentclass[compress,table]{beamer} % 在Mac下使用
% package for font
\usepackage{fontspec}
\defaultfontfeatures{Mapping=tex-text}  %%如果没有它，会有一些 tex 特殊字符无法正常使用，比如连字符。
\usepackage{xunicode,xltxtra}
\usepackage[BoldFont,SlantFont,CJKnumber,CJKchecksingle]{xeCJK}  % \CJKnumber{12345}: 一万二千三百四十五
\usepackage{CJKfntef}  %%实现对汉字加点、下划线等。
\usepackage{pifont}  % \ding{}
% package for math
\usepackage{amsfonts}

% package for graphics
\usepackage[americaninductors,europeanresistors]{circuitikz}
\usepackage{tikz}
\usetikzlibrary{plotmarks}  % placements=positioning
\usepackage{graphicx}  % \includegraphics[]{}
\usepackage{subfigure}  %%图形或表格并排排列
% package for table
\usepackage{colortbl,dcolumn}  %% 彩色表格
\usepackage{multirow}
\usepackage{multicol}
\usepackage{booktabs}
% package for code
\usepackage{fancyvrb}
\usepackage{listings}

% \usepackage{animate}
% \usepackage{movie15}

%%%%%
% setting for beamer
\usetheme{default} % Madrid（常用）, Copenhagen, AnnArbor, boxes（白色）, Frankfurt，Berkeley
\useoutertheme[subsection=true]{miniframes} % 使用Berkeley时注释本行
\usecolortheme{sidebartab}
\usefonttheme{serif}  %%英文使用衬线字体
% \setbeamertemplate{background canvas}[vertical
% shading][bottom=white,top=structure.fg!7] %%背景色，上25%的蓝，过渡到下白。
\setbeamertemplate{theorems}[numbered]
\setbeamertemplate{navigation symbols}{}  %% 去掉页面下方默认的导航条
\setbeamercovered{transparent}  %设置 beamer 覆盖效果

% 设置标题title背景色
% \setbeamercolor{title}{fg=black, bg=lightgray!60!white}
\setbeamercolor{title}{fg=white, bg=black!70!white}

% 设置每页小LOGO
\pgfdeclareimage[width=1cm]{ouc}{figures/static/ouc.pdf}
\logo{\pgfuseimage{ouc}{\vspace{-20pt}}}

% setting for font
%%\setCJKmainfont{Adobe Kaiti Std}
\setCJKmainfont{SimSun} 
%% \setCJKmainfont{FangSong_GB2312} 
%% \setmainfont{Apple Garamond}  %%苹果字体没有SmallCaps
\setCJKmainfont{SimSun} 
%FUNNY%\setCJKmainfont{DFPShaoNvW5-GB}  %%华康少女文字W5(P)
%FUNNY%\setCJKmainfont{FZJingLeiS-R-GB}  %%方正静蕾体
%FUNNY%\setmainfont{Purisa}
%\setsansfont[Mapping=tex-text]{Adobe Song Std}
     %如果装了Adobe Acrobat，可在font.conf中配置Adobe字体的路径以使用其中文字体。
     %也可直接使用系统中的中文字体如SimSun、SimHei、微软雅黑等。
     %原来beamer用的字体是sans family；注意Mapping的大小写，不能写错。
     %设置字体时也可以直接用字体名，以下三种方式等同：
     %\setromanfont[BoldFont={黑体}]{宋体}
     %\setromanfont[BoldFont={SimHei}]{SimSun}
     %\setromanfont[BoldFont={"[simhei.ttf]"}]{"[simsun.ttc]"}
% setting for graphics
\graphicspath{{figures/}}  %%图片路径
\renewcommand\figurename{图}

% setting for pdf
\hypersetup{% pdfpagemode=FullScreen,%
            pdfauthor={Xiaodong Wang},%
            pdftitle={Title},%
            CJKbookmarks=true,%
            bookmarksnumbered=true,%
            bookmarksopen=false,%
            plainpages=false,%
            colorlinks=true,%
            citecolor=green,%
            filecolor=magenta,%
            linkcolor=blue,%red(default)
            urlcolor=cyan}

% setting for fontspec
\XeTeXlinebreaklocale "zh"  %%表示用中文的断行
\XeTeXlinebreakskip = 0pt plus 1pt minus 0.1pt  %%多一点调整的空间
%%%%%

% font setting by xeCJK
\setCJKfamilyfont{NSimSun}{NSimSun}
\newcommand{\song}{\CJKfamily{NSimSun}}
%%%\setCJKfamilyfont{AdobeSongStd}{Adobe Song Std}
%%%\newcommand{\AdobeSong}{\CJKfamily{AdobeSongStd}}
\setCJKfamilyfont{FangSong}{FangSong_GB2312}
\newcommand{\fang}{\CJKfamily{FangSong}}
%%%\setCJKfamilyfont{AdobeFangsongStd}{Adobe Fangsong Std}
%%%\newcommand{\AdobeFang}{\CJKfamily{AdobeFangsongStd}}
\setCJKfamilyfont{SimHei}{SimHei}
\newcommand{\hei}{\CJKfamily{SimHei}}
%%%\setCJKfamilyfont{AdobeHeitiStd}{Adobe Heiti Std}
%%%\newcommand{\AdobeHei}{\CJKfamily{AdobeHeitiStd}}
\setCJKfamilyfont{KaiTi}{KaiTi}
\newcommand{\kai}{\CJKfamily{KaiTi}}
%%%\setCJKfamilyfont{AdobeKaitiStd}{Adobe Kaiti Std}
\newcommand{\AdobeKai}{\CJKfamily{AdobeKaitiStd}}
\setCJKfamilyfont{LiSu}{LiSu}
\newcommand{\li}{\CJKfamily{LiSu}}
\setCJKfamilyfont{YouYuan}{YouYuan}
\newcommand{\you}{\CJKfamily{YouYuan}}
\setCJKfamilyfont{FZJingLei}{FZJingLeiS-R-GB}
\newcommand{\jinglei}{\CJKfamily{FZJingLei}}
\setCJKfamilyfont{MSYH}{Microsoft YaHei}
\newcommand{\msyh}{\CJKfamily{MSYH}}

% 自定义颜色
\def\Red{\color{red}}
\def\Green{\color{green}}
\def\Blue{\color{blue}}
\def\Mage{\color{magenta}}
\def\Cyan{\color{cyan}}
\def\Brown{\color{brown}}
\def\White{\color{white}}
\def\Black{\color{black}}

\lstnewenvironment{xmlCode}[1][]{% for Java
  \lstset{
    basicstyle=\tiny\ttfamily,%
    columns=flexible,%
    framexleftmargin=.7mm, %
    % frame=shadowbox,%
    % rulesepcolor=\color{cyan},%
     frame=single,%
    backgroundcolor=\color{white},%
    xleftmargin=4\fboxsep,%
    xrightmargin=4\fboxsep,%
    numbers=left,numberstyle=\tiny,%
    numberblanklines=false,numbersep=7pt,%
    language=xml, %
    }\lstset{#1}}{}

\lstnewenvironment{javaCode}[1][]{% for Java
  \lstset{
    basicstyle=\tiny\ttfamily,%
    columns=flexible,%
    framexleftmargin=.7mm, %
    frame=shadowbox,%
    rulesepcolor=\color{cyan},%
    % frame=single,%
    backgroundcolor=\color{white},%
    xleftmargin=4\fboxsep,%
    xrightmargin=4\fboxsep,%
    numbers=left,numberstyle=\tiny,%
    numberblanklines=false,numbersep=7pt,%
    language=Java, %
    }\lstset{#1}}{}

\lstnewenvironment{shCode}[1][]{% for Java
  \lstset{
    basicstyle=\scriptsize\ttfamily,%
    columns=flexible,%
    framexleftmargin=.7mm, %
    frame=shadowbox,%
    rulesepcolor=\color{brown},%
    % frame=single,%
    backgroundcolor=\color{white},%
    xleftmargin=4\fboxsep,%
    xrightmargin=4\fboxsep,%
    numbers=left,numberstyle=\tiny,%
    numberblanklines=false,numbersep=7pt,%
    language=sh, %
    }\lstset{#1}}{}

\newcommand\ask[1]{\vskip 4bp \tikz \node[rectangle,rounded corners,minimum size=6mm,
  fill=white,]{\Cyan \includegraphics[height=1.5cm]{question} \Large \msyh #1};}

\newcommand\wxd[1]{\vskip 4bp \tikz \node[rectangle,minimum size=6mm,
  fill=blue!60!white,]{\White \ding{118} \msyh #1};}

\newcommand\xyy[1]{\vskip 2bp \tikz \node[rectangle,minimum size=3mm,
  fill=black!80!white,]{\White \msyh\scriptsize #1};}

\newcommand\homework[1]{\vskip 2bp \tikz \node[rectangle,minimum size=3mm,
  fill=red!80!white,]{\White \ding{45} \msyh\scriptsize 课后小作业 } ; {\kai\small #1}} 

\newcommand\cxf[1]{\vskip 4bp \tikz \node[rectangle,rounded corners,minimum size=6mm,
  fill=purple!60!white,]{\White \ding{42} \msyh #1};}

\newcommand\tta[1]{\vskip 4bp \tikz \node[rectangle,minimum size=6mm,
  fill=blue!60!white,]{\White \ding{118} \msyh #1};}

\newcommand\ttb[1]{\vskip 4bp \tikz \node[rectangle,rounded corners,minimum size=6mm,
  fill=purple!60!white,]{\White \ding{42} \msyh #1};}

\newcommand\ttc[1]{\vskip 2bp \tikz \node[rectangle,minimum size=3mm,
  fill=black!80!white,]{\White \msyh\scriptsize #1};}

\newcommand\notice[1]{\vskip 4bp \tikz \node[rectangle,rounded corners,minimum size=6mm,
  fill=red!80!white,]{\White \scriptsize \ding{42} \msyh #1};}

\newcommand\samp[1]{\vskip 2bp \tikz \node[rectangle,minimum size=3mm,
  fill=white!100!white,]{\Mage\msyh \small CODE \ding{231} \Black #1};\vskip -8bp}

\newcommand\codeset[1]{\vskip 2bp \tikz \node[rectangle,minimum size=3mm,
  fill=white!100!white,]{\Mage\msyh \small 课程配套代码 \ding{231} \Black #1};\vskip -8bp}

\newcommand\pptlink[2]{\vskip 4bp \tikz \node[rectangle,rounded corners,minimum size=6mm,
  fill=blue!70!white,]{\href{run:#1}{\White \scriptsize \msyh 动画演示 #2}};}



\setbeamerfont{frametitle}{series=\msyh} % 修改Beamer标题字体

\makeatletter
\newcommand{\Extend}[5]{\ext@arrow 0099{\arrowfill@#1#2#3}{#4}{#5}}
\makeatother


%%%%%%%%%%%%%%%%%%%%%%%%%%%%%%%%%%%%%%%%%%%%%%%%%%%%%%%%%%%%%%%%%%%%%%%%%%%%%%%
\title[JAVA]{\hei {\huge Java 应用与开发}\\  
 面向对象编程进阶 PART2}
\author[王晓东]{王晓东\\
  \href{mailto:wangxiaodong@ouc.edu.cn}{\footnotesize wangxiaodong@ouc.edu.cn}}
\institute[中国海洋大学]{\small 中国海洋大学}
\date{\today}
\titlegraphic{\vspace{-6em}\includegraphics[height=6cm]{static/ouc.pdf}\vspace{-6em}}
%%%%%%%%%%%%%%%%%%%%%%%%%%%%%%%%%%%%%%%%%%%%%%%%%%%%%%%%%%%%%%%%%%%%%%%%%%%%%%%
\begin{document}
%% Delete this, if you do not want the table of contents to pop up at
%% the beginning of each subsection:
\AtBeginSection[]{                              % 在每个Section前都会加入的Frame
  \frame<handout:0>{
    \frametitle{\textbf{\hei 接下来…}}
    \tableofcontents[currentsection]
  }
}  %

\AtBeginSubsection[]                            % 在每个子段落之前
{
  \frame<handout:0>                             % handout:0 表示只在手稿中出现
  {
    \frametitle{\textit{\hei 接下来…}}\small
    \tableofcontents[current,currentsubsection] % 显示在目录中加亮的当前章节
  }
}
 \frame{\titlepage}

\begin{frame}
 \frametitle{学习目标}
 \begin{enumerate}
 \item 理解多态和虚方法调用的概念，掌握其用法
 \item 掌握方法重载的方法
 \item 掌握static属性、方法和初始化块的用法
 \item 了解设计模式，掌握单例设计模式
 \item 掌握final关键字的概念和使用方法
 \end{enumerate}  
\end{frame}

\section*{大纲}
\frame{\frametitle{大纲} \tableofcontents }

\section{多态性}

\begin{frame}[fragile] % [fragile]参数使得能够插入代码
  \frametitle{什么是多态？}

  在Java中，子类的对象可以替代父类的对象使用称为{\hei 多态}。

  \wxd{Java引用变量与所引用对象间的类型匹配关系}
  
  \begin{itemize}[<+-| structure@+>]\kai
  \item 一个对象只能属于一种确定的数据类型，该类型自对象创建直至销毁不能改变。
  \item 一个引用类型变量可能引用（指向）多种不同类型的对象{\Red ——既可以引用其声明类型的对象，也
    可以引用其声明类型的子类的对象。}
    \begin{javaCode}
      Person p = new Student(); //Student 是 Person 的子类  
    \end{javaCode}
  \end{itemize}

  \begin{figure}
    \centering
    \includegraphics[width=0.5\textwidth]{fig-poly.pdf}
  \end{figure}
\end{frame}

\begin{frame}[fragile] % [fragile]参数使得能够插入代码
\frametitle{多态性}

\xyy{多态性同样适用与引用类型数组元素}

\begin{javaCode}
  Person[] p = new Person[3]; 
  p[0] = new Student(); // 假设 Student 类继承了 Person 类
  p[1] = new Person();
  p[2] = new Graduate(); //假设 Graduate 类继承了 Student 类
\end{javaCode}

\pause

\notice{父类引用对子类对象的能力屏蔽性}

一个引用类型变量如果声明为父类的类型，但实际引用的是子类对象，该变量则不能再访问子类
中添加的属性和方法。

\begin{javaCode}
  Student m = new Student();
  m.setSchool("ouc"); // 合法
  Person e = new Student();
  e.setSchool("ouc"); // 非法
\end{javaCode}

\end{frame}

\begin{frame}[fragile] % [fragile]参数使得能够插入代码
\frametitle{多态用法示例}

\samp{Person.java}

\begin{javaCode}
  public class Person {...}
\end{javaCode}

\samp{Student.java}

\begin{javaCode}
  public class Student extends Person {
    private String school;

    public void setSchool(String school) {
      this.school = school;
    }

    public String getSchool() {
      return school;
    }

    @Override
    public String getInfo() {
      return super.getInfo() + "\tSchool: " + school;
    }
  }
\end{javaCode}
\end{frame}

\begin{frame}[fragile] % [fragile]参数使得能够插入代码
\frametitle{多态用法示例}

\samp{PolymorphismSample.java}

\begin{javaCode}

  public class PolymorphismSample {
    public void show(Person p) {
      System.out.println(p.getInfo());
    }
    
    public static void main(String[] args) {
      PolymorphismSample ps = new PolymorphismSample();
      Person p = new Person();
      ps.show(p);
      Student s = new Student();
      ps.show(s);
    }
  }
\end{javaCode}

\pause

\notice{多态提升方法通用性}

{\kai\Mage show()方法既可以处理Person类型的数据，又可以处理Student类型的数
  据，乃至未来定义的任何Person子类类型的数据，即不必为相关的每种类型单
  独声明一个处理方法，提高了代码的通用性。}

\end{frame}

\begin{frame}[fragile] % [fragile]参数使得能够插入代码
\frametitle{虚方法调用}
\cxf{思考}

{\Mage 一个引用类型的变量如果声明为父类的类型，但实际引用的是子类对象，
  则该变量就不能再访问子类中添加的属性和方法。}

\pause

\begin{block}{}
  但如果此时调用的是父类中声明过、且在子类中重写过的方法，情况如何？  
\end{block}

\end{frame}

%%%\begin{frame}[fragile] % [fragile]参数使得能够插入代码
%%%\frametitle{虚方法调用示例}
%%%
%%%\samp{Book.java}
%%%
%%%\begin{javaCode}
%%%  public class Book {
%%%    private String name;
%%%    private double price;
%%%
%%%    public void setName(String name) {
%%%      this.name = name;
%%%    }
%%%    public String getName() {
%%%      return name;
%%%    }
%%%    public void setPrice(double price) {
%%%      this.price = name;
%%%    }
%%%    public double getPrice() {
%%%      return price;
%%%    }
%%%    public void show() {
%%%      System.out.println("Book name:" + name + "\nprice" + price);
%%%    }
%%%  }
%%%\end{javaCode}
%%%\end{frame}
%%%
%%%\begin{frame}[fragile] % [fragile]参数使得能够插入代码
%%%\frametitle{虚方法调用示例}
%%%
%%%\samp{Novel.java}
%%%
%%%\begin{javaCode}
%%%  public class Novel extends Book {
%%%    private String author;
%%%
%%%    public void setAuthor(String author) {
%%%      this.author = author;
%%%    }
%%%    public String getAuthor() {
%%%      return author;
%%%    }
%%%    public void show() {
%%%      super.show();
%%%      System.out.println("Author:" + author);
%%%    }
%%%  }
%%%\end{javaCode}
%%%\end{frame}
%%%
%%%\begin{frame}[fragile] % [fragile]参数使得能够插入代码
%%%  \frametitle{虚方法调用示例}
%%%
%%%  \samp{Test.java}
%%%
%%%  \begin{javaCode}
%%%    public void class Test {
%%%      public static void main(String[] args) {
%%%        Test t = new Test();
%%%        Book b = new Book();
%%%        b.setName("English Language");
%%%        b.setPrice(43);
%%%        t.process(b);
%%%        Novel v = new Novel();
%%%        v.setName("Great Expectations");
%%%        v.setPrice(35.48);
%%%        v.setAuthor("Charles Dickens");
%%%        t.process(v);
%%%      }
%%%      public void process(Book b) {
%%%        b.show();
%%%      }
%%%    }
%%%  \end{javaCode}
%%%\end{frame}
%%%
%%%\begin{frame}[fragile] % [fragile]参数使得能够插入代码
%%%\frametitle{虚方法调用示例}
%%%程序输出结果：
%%%
%%%\begin{stdoutCode}
%%%BookName:English Language
%%%Price:43.0
%%%BookName:Great Eespectations
%%%Price:35.48
%%%Author:Charles Dickens
%%%\end{stdoutCode}
%%%
%%%\pause
%%%\notice{运行结果分析}
%%%
%%%系统根据运行时对象的真正类型来确定具体调用哪一个方法，这一机制被称为
%%%{\hei\Blue 虚方法调用（Virtual Method Invocation）}。
%%%\end{frame}

\begin{frame}[fragile] % [fragile]参数使得能够插入代码
\frametitle{对象造型}

引用类型数据值之间的强制类型转换称为{\hei\Mage 造型（Casting）}。

\begin{enumerate}[<+-| structure@+>]
\item 从子类到父类的类型转换可以自动进行。
  \begin{javaCode}
    Person p = new Student();    
  \end{javaCode}
\item 在多态的情况下，从父类到子类的类型转换必须通过造型实
  现。\footnote{{\kai 有时我们可能需要恢复一个对象的本来面目，以发挥其
      全部潜力。}}
  \begin{javaCode}
    Person p1 = new Student();  
    Student s1 = (Student)p1;   // 合法
    Person p2 = new Person();   
    Student s2 = (Student)p2;  // 非法  
  \end{javaCode}
\item 无继承关系的引用类型间的转换是非法的。
  \begin{javaCode}
    String s = "Hello World!";
    Person p = (Person)s; // 非法  
  \end{javaCode}
\end{enumerate}
\end{frame}

\begin{frame}[fragile] % [fragile]参数使得能够插入代码
\frametitle{instanceof运算符}

{\footnotesize 如果运算符instanceof左侧的变量当前时刻所引用的对象的{\hei\Red 真正类型}是其右侧给出的
  类型{\Blue 或者是其子类}，则整个表达式的结果为true。}

\pause

\begin{javaCode}
class Person { --- }
class Student extends Person { --- }

public class Tool {
  public void distribute(Person p) {
    if (p instanceof Student) {
      System.out.println("处理 Student 类型及其子类类型对象");
    } else {
      System.out.println("处理 Person 类型及其子类类型对象");
    }
  }
}
\end{javaCode}

\begin{javaCode}
public class Test() {
  public static void main(String[] args) {
    Tool t = new Tool();
    Student s = new Student();
    t.distribute(t);
  }
}
\end{javaCode}
\end{frame}

\begin{frame}[fragile]
  \frametitle{虚方法调用和造型}

  \codeset{package sample.oop.poly}

  \begin{itemize}\small
  \item VirtualMethodSample.java
  \item Person.java
  \item Student.java
  \end{itemize}
  
  \wxd{虚方法调用和造型强化}
  
  \begin{itemize}[<+-| alert@+>]
  \item 系统根据运行时对象的真正类型来确定具体调用哪一个方法，这一机制被称为
    {\hei\Blue 虚方法调用}。
  \item 造型是引用类型数据值之间的强制类型转换。
  \item instanceof运算符判断的是当前所引用对象的真正类型是什么，而不是声明的引用类型。
  \end{itemize}
\end{frame}

%\begin{frame}[fragile] % [fragile]参数使得能够插入代码
%\frametitle{协变返回类型}
%
%{\hei 协变返回类型}——从Java SE5.0开始引入，允许重写方法时修改其返回值的类型，但必须是重写
%前方法返回值类型的{\hei 子类或实现类类型\footnote{实现类后续课程讲述。}}。
%\begin{javaCode}
%class A {
%  public Person getAssistor() {
%    Person p = new Person();
%    //---
%    return p;
%  }
%}
%class B extends A {
%  public Student getAssistor() { //重写方法时改变了返回值类型
%    Student s = new Student();
%    s.setName("Nancy");
%    s.setAge(18);
%    s.setSchool("THU");
%    return s;
%  }
%}
%\end{javaCode}
%\end{frame}
%
%\begin{frame}[fragile] % [fragile]参数使得能够插入代码
%\frametitle{协变返回类型}
%\samp{如何不采用协变返回类型实现上述效果}
%\begin{javaCode}
%class B extends A {
%  public Person getAssistor() {
%    Student s = new Student();
%    s.setName("Nancy");
%    s.setAge(18);
%    s.setSchool("THU");
%    return s; //合法，基于多态机制返回子类的实例
%  }
%}  
%public class TestCovarianReturnType {
%  public static void main(String[] args) {
%    B b = new B();
%    Student stu = (Student)b.getAssistor();
%    System.out.println(stu.getSchool());
%  }
%}
%\end{javaCode}
%当重写方法中实际返回的是原返回类型的子类类型数据时，使用协变返回模式{\hei\Red 不需要再进行强制类型
%转换就可以访问其子类中添加的成员}。
%\end{frame}

\section{方法重载}
\begin{frame}[fragile] % [fragile]参数使得能够插入代码
\frametitle{什么是方法重载}

在一个类中存在多个同名方法的情况称为{\hei 方法重载（Overload）}。

\begin{itemize}\kai
\item 重载方法参数列表必须不同。
\item 重载既可以用于普通方法，也可以用于构造方法。
\end{itemize}

\codeset{sample.oop.MethodOverloadSample.java}
\end{frame}

\begin{frame}[fragile] % [fragile]参数使得能够插入代码
  \frametitle{如何调用重载的构造方法}
  
  \wxd{使用this调用当前类中重载构造方法}

  可以在构造方法的第一行使用关键字this调用其它（重载的）构造方法。

  \begin{javaCode}
    public class Person {
      ...
      public Person(String name,int age) {
        this.name = name;
        this.age = age;
      }
      public Person(String name) {
        this(name, 18);
      }
      ...
    }
  \end{javaCode}

  \notice{注意}
  
  关键字this的此种用法只能用在构造方法中，且this()语句如果出现必须位于
  方法体中代码的第一行。
\end{frame}

\begin{frame}[fragile] % [fragile]参数使得能够插入代码
  \frametitle{如何调用重载的构造方法}
  \wxd{使用super调用父类构造方法}

  \samp{Person.java}

  \begin{javaCode}
    public class Person {
      ... (此处没有无参构造方法)
      public Person(String name, int age) {
        this.name = name;
        this.age = age;
      }
      ...
    }  
  \end{javaCode}
  
  \samp{Student.java}
  
  \begin{javaCode}
    public class Student extends Person {
      private String school;
      public Student(String name, int age, String school) {
        super(name, age); // 显式调用父类有参构造方法
        this.school = school;
      }
      public Student(String school) { //编译出错
        // super(); // 隐式调用父类有参构造方法，则自动调用父类无参构造方法
        this.school = school;
      }
    }
\end{javaCode}
\end{frame}


\begin{frame}[fragile] % [fragile]参数使得能够插入代码
  \frametitle{上述代码为什么会编译出错}

  {\Blue\hei 在Java类的构造方法中一定直接或间接地调用了其父类的构造方法（Object类除外）。}

  \begin{enumerate}[<+-| structure@+>]
  \item 在子类的构造方法中可使用super语句调用父类的构造方法，其格式为super(<实参列表>)。
  \item 如果子类的构造方法中既没有显式地调用父类构造方法，也没有使用this关键字调用同一个类
    的其他重载构造方法，则系统会默认调用父类无参数的构造方法，其格式为super()。
  \item 如果子类构造方法中既未显式调用父类构造方法，而父类中又没有无参的构造方法，则编译出
    错。
  \end{enumerate}

  \codeset{sample.oop.ConstructorOverloadSample.java}
  
\end{frame}

%%%\begin{frame}[fragile] % [fragile]参数使得能够插入代码
%%%\frametitle{对象构造/初始化细节}
%%%
%%%\xyy{第一阶段} 为新建对象的实例变量分配存储空间并进行默认初始化。
%%%\xyy{第二阶段} 按下述步骤继续初始化实例变量：
%%%
%%%\begin{enumerate}
%%%\item 绑定构造方法参数；
%%%\item 如有this()调用，则调用相应的重载构造方法然后跳转到步骤5；
%%%\item 显式或隐式追溯调用父类的构造方法（Object类除外）；
%%%\item 进行实例变量的显式初始化操作；
%%%\item 执行当前构造方法的方法体中其余的语句。 
%%%\end{enumerate}
%%%\end{frame}

\section{关键字static}

\begin{frame}[fragile] % [fragile]参数使得能够插入代码
  \frametitle{关键字static}

  在Java类中声明{\hei\Red 属性、方法和内部类}时，可使用关键字static作为修饰符。

  \begin{itemize}[<+-| structure@+>]
  \item static标记的属性或方法由整个类（所有实例）共享，如访问控制权限
    允许，可不必创建该类对象而直接用类名加“.”调用。
  \item static成员也称{\hei\Blue “类成员”或“静态成员”}，如“类属
    性”、“类变量”、“类方法”和“静态方法”等。
  \end{itemize}
\end{frame}

\begin{frame}[fragile] % [fragile]参数使得能够插入代码
  \frametitle{static属性和方法}
  \wxd{static属性}
  
  \begin{itemize}
  \item static属性由其所在类（包括该类所有的实例）共享。
  \item 非static属性则必须依赖具体/特定的对象（实例）而存在。
  \end{itemize}

  \wxd{static方法}

  要在static方法中调用其所在类的非static成员，应首先创建一个该类的对象，
  通过该对象来访问其非static成员。

  \codeset{sample.oop.StaticMemberAndMethodSample.java}
\end{frame}

\begin{frame}[fragile] % [fragile]参数使得能够插入代码
  \frametitle{初始化块}

  \wxd{static初始化块}

  在类的定义体中，方法的外部可包含static语句块，{\Red static块仅在其所
    属的类被载入时执行一次}，通常用于初始化化static（类）属性。

  \wxd{非static初始化块}
  
  非static的初始化块在创建对象时被自动调用。

  \codeset{sample.oop.StaticInitBlockSample.java}
  
\end{frame}

\begin{frame}[fragile] % [fragile]参数使得能够插入代码
\frametitle{静态导入}

静态导入用于在一个类中导入其他类或接口中的static成员，语法格式：\\
import static <包路径>.<类名>.*\\
或：\\
import static <包路径>.<类名>.<静态成员名>\\

\samp{应用示例}
\begin{javaCode}
import static java.lang.Math.*;
public class Test {
  public static void main(String[] args) {
    double d = sin(PI * 0.45);
    System.out.println(d);
  }
}
\end{javaCode}
\end{frame}

\begin{frame}[fragile] % [fragile]参数使得能够插入代码
  \frametitle{Singleton设计模式\footnote{Singleton设计模式也称“单子模
      式”或“单例模式”。}}

\wxd{所谓“模式”就是被验证为有效的常规问题的典型解决方案。}

\begin{itemize}
\item {\hei 设计模式（Design Pattern）}在面向对象分析设计和软件开发中占有重要地位。
\item 好的设计模式可以使我们更加方便的重用已有的成功设计和体系结构，极大的提高代码的重用性和可维护性。
\end{itemize}

\end{frame}

\begin{frame}[fragile]
  \frametitle{经典设计模式分类}
  \xyy{创建型模式} 涉及对象的实例化，特点是不让用户代码依赖于对象的创建或排列方式，避免用户直接使用new创建对象。\\
  {\Red\kai 工厂方法模式、抽象工厂方法模式、生成器模式、原型模式和单例模式}\pause
  
  \xyy{行为型模式} 涉及怎样合理的设计对象之间的交互通信，以及合理为对象分配职责，让设计富有弹性、易维护、易复用。\\
  {\Blue\kai 责任链模式、命令模式、解释器模式、迭代器模式、中介者模式、备忘
    录模式、观察者模式、状态模式、策略模式、模板方法模式和访问者模式}\pause

  \xyy{结构型模式} 涉及如何组合类和对象以形成更大的结构，和类有关的结构型模式涉及如何合理使用继承机制，和对象有关的结构型模式涉及如何合理的使用对象组合机制。\\
  {\Mage\kai 适配器模式、组合模式、代理模式、享元模式、外观模式、桥接模式和
    装饰模式}
\end{frame}


\begin{frame}[fragile] % [fragile]参数使得能够插入代码
  \frametitle{Singleton设计模式}

  \notice{采用调试方式讲解示例代码}
  
  \codeset{package sample.oop.singleton}


%%%  \begin{javaCode}
%%%    public class Single {
%%%      private String name;
%%%      private static Single onlyone = new Single();
%%%      private Single() {}
%%%      public void setName(String name) {
%%%        this.name = name;
%%%      }
%%%      public String getName(){
%%%        return name;
%%%      }
%%%      public static Single getSingle() {
%%%        return onlyone;
%%%      }
%%%    }
%%%  \end{javaCode}
%%%
%%%  \begin{javaCode}
%%%    public class TestSingle {
%%%      public static void m1() {
%%%        Single s2 = Single.getSingle(); 
%%%        System.out.println(s2.getName());
%%%      }
%%%      public static void main(String args[]) {
%%%        Single s1 = Single.getSingle();
%%%        s1.setName("Tom");
%%%        m1();
%%%      }
%%%    }
%%%  \end{javaCode}
%%%  \only<2>{\Mage \small 程序输出结果：Tom}
\end{frame}

\begin{frame}[fragile] % [fragile]参数使得能够插入代码
  \frametitle{Singleton设计模式}

  \wxd{Singleton代码的特点}

  \begin{enumerate}[<+-| alert@+>]
  \item 使用静态属性onlyone来引用一个“全局性”的Single实例。
  \item 将构造方法设置为private的，这样在外界将不能再使用new关键字来创建该类的新实例。
  \item 提供public static的方法getSingle()以使外界能够获取该类的实例，达到全局可见的效果。
  \end{enumerate}

  \pause
  \wxd{Singleton的使用场景}

  {\kai 在任何使用到Single类的Java程序中（这里指的是一次运行中），需要
    确保只有一个Single类的实例存在（如Web应用ServletContext全局上下文对
    象），则使用该模式。}


\end{frame}

\section{关键字final}
\begin{frame}[fragile] % [fragile]参数使得能够插入代码
\frametitle{关键字final}
在声明Java类、变量和方法时可以使用关键字final来修饰，以使其具有“终态”的特性：

\begin{enumerate}\kai
\item final标记的类不能被继承；
\item final标记的方法不能被子类重写；
\item final标记的变量（成员变量或局部变量）即成为常量,只能赋值一次；
\item final标记的成员变量必须在声明的同时或在每个构造方法中显式赋值，然后才能使用；
\item final不允许用于修饰构造方法、抽象类以及抽象方法。
\end{enumerate}
\end{frame}

\begin{frame}[fragile] % [fragile]参数使得能够插入代码
\frametitle{关键字final应用举例}
\begin{javaCode}
public final class Test {
  public static int totalNumber = 5;
  public final int id;
  public Test() {
    id = ++totalNumber; // 赋值一次
  }
  public static void main(String[] args) {
    Test t = new Test();
    System.out.println(t.id);
    final int i = 10;
    final int j;
    j = 20;
    j = 30;  // 非法
  }
}
\end{javaCode}
\end{frame}
%%%%%%%%%%%%%%%%%%%%%%%%%%%%%%%%%%%%%%%%%%%%%%%%%%%
\begin{frame}
  \frametitle{本节习题}

  \wxd{简答题}
  \begin{enumerate}
  \item 为什么建议Java类需要编写无参构造方法，哪怕该方法什么都没做？
  \item 关键字static都可以用来修饰Java类的那些成员？
  \end{enumerate}

  \wxd{小编程}
  \begin{enumerate}
  \item 练习本节中所有示例代码，理解并掌握其用法。
  \item 自行调试单例设计模式程序，学习Eclipse的Debug方法，练习断点、单
    步执行、跟进方法等操作，查看内存变化。
  \end{enumerate}
\end{frame}
%%%%%%%%%%%%%%%%%%%%%%%%%%%%%%%%%%%%%%%%%%%%%%%%%%%%%%%%%%%%%%%%%%%%%%%%%%%%%%%
% TKS Page %%%%%%%%%%%%%%%%%%%%%%%%%%%%%%%%%%%%%%%%%%%%
\begin{frame}
\centering
{\Huge \textcolor{blue}{THE END}} \\
\vspace{5mm}
{\Large wangxiaodong@ouc.edu.cn} \\
\end{frame}
%%%%%%%%%%%%%%%%%%%%%%%%%%%%%%%%%%%%%%%%%%%%%%%%%%%%%%%
%%%%%%%%%%%%%%%%%%%%%%%%%%%%%%%%%%%%%%%%%%%%%%%%%%%%%%%%%%%%%%%%%%%%%%%%%%%%%%%
\end{document}
